% ================= CAPÍTULO 5: LÍMITES =================
\chapter{Límites}
\prob{SEMANAS 5-6}{tu cuaderno llega hasta el límite de la suma; el resto está agregado}

\section{Definición}
\apuntes
\begin{sello}{DEFINICIÓN ε-δ}
$f:I\to\R$, $I$ intervalo, $x_0\in\bar I$ (en $I$ o en su borde).
\[\lim_{x\to x_0}f(x)=L\iff\forall\eps>0\ \exists\delta>0:\]\[0<|x-x_0|<\delta,\ x\in I\ \Longrightarrow\ |f(x)-L|<\eps.\]
\end{sello}
\begin{center}
\begin{tikzpicture}[x=1.7cm,y=1.1cm]
\fill[rojo!9] (0,1.55) rectangle (2.7,2.45);
\fill[teal!18] (0.8,0) rectangle (1.2,2.9);
\draw[rojo,dashed] (0,1.55) -- (2.7,1.55) (0,2.45) -- (2.7,2.45);
\draw[teal,dashed] (0.8,0) -- (0.8,2.9) (1.2,0) -- (1.2,2.9);
\draw[-{Stealth}] (0,0) -- (2.8,0); \draw[-{Stealth}] (0,0) -- (0,3);
\draw[acero,very thick,domain=0.1:2.3,samples=40] plot (\x,{2+0.9*(\x-1)+0.25*(\x-1)^3});
\draw[acero,thick,fill=papel] (1,2) circle (2pt);
\draw[densely dotted] (1,0) -- (1,2) -- (0,2);
\node[font=\scriptsize,left] at (0,2) {$L$}; \node[font=\scriptsize,below] at (1,0) {$x_0$};
\node[rojo,font=\ttfamily\scriptsize,anchor=west] at (1.6,2.62) {banda $(L-\eps,L+\eps)$};
\node[teal,font=\ttfamily\scriptsize,anchor=west] at (1.25,0.3) {entorno $(x_0-\delta,x_0+\delta)$};
\end{tikzpicture}\\
{\ttfamily\scriptsize\color{suave} Te dan la banda roja; vos tenés que encontrar la franja teal tal que la curva sobre ella quede dentro de la banda. El punto $x_0$ mismo no importa (círculo hueco).}
\end{center}
\begin{memo}{TRES COSAS DE LA DEFINICIÓN}
\textbf{El orden:} primero $\eps$ (lo que te exigen), después $\delta$ (lo que respondés). El $\delta$ puede depender de $\eps$.\\
\textbf{El $0<$:} $|x-x_0|>0$ saca al punto: el límite no mira $f(x_0)$, ni siquiera necesita que exista.\\
\textbf{Achicar sirve:} si un $\delta$ funciona, cualquiera más chico también.
\end{memo}

\section{Ejemplos de clase}
\apuntes
\begin{sello}{CINCO EJEMPLOS}
1) \textbf{Constante} $f=k$: $|f(x)-k|=0<\eps$ siempre; sirve cualquier $\delta$ (ej. $\delta=1$). Límite $k$.\\[2pt]
2) \textbf{Identidad} $f=x$: $|x-x_0|<\eps$ si $\delta=\eps$. Límite $x_0$.\\[2pt]
3) $f=-2x$: $|f(x)-(-2x_0)|=2|x-x_0|<\eps$ si $\delta=\frac\eps2$. Límite $-2x_0$.\\[2pt]
4) $f(x)=x$ si $x\neq0$, $f(0)=1$: límite en 0 es 0, con $\delta=\eps$. (El valor $f(0)=1$ no importa.)\\[2pt]
5) \textbf{Signo} ($1$ si $x>0$, $0$ en 0, $-1$ si $x<0$): \textbf{no existe} el límite en 0.
\end{sello}
\agregado
\begin{demo}
Ejemplo 5. Supongo que el límite es $L$ y tomo $\eps=1$. Para cualquier $\delta$, en $(0,\delta)$ la función vale 1 y en $(-\delta,0)$ vale $-1$. Tendría que ser $|1-L|<1$ y $|-1-L|<1$, o sea $L\in(0,2)$ y $L\in(-2,0)$ a la vez. Imposible. $\blacksquare$
\end{demo}
\begin{tip}{EL PATRÓN LINEAL}
Para $f(x)=mx+n$: $|f(x)-f(x_0)|=|m||x-x_0|$, así que $\delta=\frac\eps{|m|}$. La pendiente divide.
\end{tip}

\section{Unicidad y límite de la suma}
\apuntes
\begin{sello}{TEOREMA: UNICIDAD DEL LÍMITE}
Si $\lim_{x\to x_0}f=L_1$ y $\lim_{x\to x_0}f=L_2$, entonces $L_1=L_2$.
\end{sello}
\begin{demo}
Por absurdo, $L_1<L_2$. Tomo $\eps=\frac{L_2-L_1}2$. Hay $\delta_1$ que deja a $f(x)$ en $(L_1-\eps,L_1+\eps)$ y $\delta_2$ que la deja en $(L_2-\eps,L_2+\eps)$. Con $\delta=\min\{\delta_1,\delta_2\}$ valen las dos a la vez, pero con ese $\eps$ los dos intervalos son \textbf{disjuntos} (se tocan justo en el punto medio). Contradicción. $\blacksquare$
\end{demo}
\apuntes
\begin{sello}{TEOREMA: LÍMITE DE LA SUMA}
$\lim_{x\to x_0}f=L_1$ y $\lim_{x\to x_0}g=L_2\Rightarrow\lim_{x\to x_0}(f+g)=L_1+L_2$.
\end{sello}
\begin{demo}
Dado $\eps$, uso $\frac\eps2$ en cada definición: hay $\delta_1$ y $\delta_2$ con $|f(x)-L_1|<\frac\eps2$ y $|g(x)-L_2|<\frac\eps2$. Con $\delta=\min\{\delta_1,\delta_2\}$ y la triangular:
\[|(f+g)(x)-(L_1+L_2)|\le|f(x)-L_1|+|g(x)-L_2|<\tfrac\eps2+\tfrac\eps2=\eps.\ \blacksquare\]
\end{demo}
\begin{trampa}{ACÁ SE CORTA TU CUADERNO}
Todo lo que sigue en este capítulo y el siguiente lo da la profesora (está en el cronograma de las semanas 5 a 7), pero no está en tus apuntes. Es teoría estándar del curso.
\end{trampa}

\section{Álgebra de límites}
\agregado
\begin{sello}{CUATRO LEMAS}
Si $f\to L$ cuando $x\to x_0$:\\[2pt]
a) \textbf{Acotación local:} existe un entorno de $x_0$ donde $|f(x)|<|L|+1$.\\
b) \textbf{Conservación del signo:} si $L>0$, existe un entorno donde $f(x)>\frac L2>0$.\\
c) \textbf{Producto:} si además $g\to M$, entonces $fg\to LM$.\\
d) \textbf{Cociente:} si además $g\to M\neq0$, entonces $\frac fg\to\frac LM$.
\end{sello}
\begin{demo}
a) Uso $\eps=1$: $|f(x)|\le|f(x)-L|+|L|<1+|L|$.\\[2pt]
b) Uso $\eps=\frac L2$: $f(x)>L-\frac L2=\frac L2$.\\[2pt]
c) Sumo y resto $f(x)M$: $|fg-LM|\le|f|\,|g-M|+|M|\,|f-L|$. Por a), $|f|<|L|+1$. Cada término se hace chico eligiendo el $\delta$ adecuado.\\[2pt]
d) Alcanza con $\frac1g\to\frac1M$. $\left|\frac1g-\frac1M\right|=\frac{|g-M|}{|g||M|}$, y por b) aplicado a $|g|$, cerca de $x_0$ vale $|g|>\frac{|M|}2$: queda $\le\frac2{M^2}|g-M|$. $\blacksquare$
\end{demo}

\agregado
\begin{sello}{COMPARACIÓN, SÁNDWICH Y ACOTADO POR INFINITÉSIMO}
\textbf{Monotonía del límite:} si $f\le g$ cerca de $x_0$, $f\to L$ y $g\to M$, entonces $L\le M$.\\[2pt]
\textbf{Sándwich:} si $g\le f\le h$ cerca de $x_0$ y $g,h\to L$, entonces $f\to L$.\\[2pt]
\textbf{Acotado por infinitésimo:} si $|g|\le K$ cerca de $x_0$ y $h\to0$, entonces $gh\to0$ (aunque $g$ no tenga límite).
\end{sello}
\begin{demo}
Monotonía: si fuera $L>M$, por conservación del signo aplicada a $f-g\to L-M>0$, sería $f>g$ cerca. Contradicción.\\[2pt]
Sándwich: dado $\eps$, cerca de $x_0$ vale $L-\eps<g\le f\le h<L+\eps$.\\[2pt]
Acotado por infinitésimo: $|gh|\le K|h|$, y $K|h|<\eps$ si $|h|<\frac\eps K$. $\blacksquare$
\end{demo}
\begin{tip}{LA IMAGEN}
Un noise de amplitud acotada ($\sin\frac1x$) por una máscara que se apaga ($x$): el resultado se apaga. Por eso $x\sin\frac1x\to0$.
\end{tip}

\agregado
\begin{sello}{LÍMITE DE LA COMPOSICIÓN}
Si $g(x)\to b$ cuando $x\to x_0$, y $f$ es continua en $b$, entonces $f(g(x))\to f(b)$.
\end{sello}
\begin{demo}
Dado $\eps$, por continuidad de $f$ en $b$ hay $\eta$ con $|y-b|<\eta\Rightarrow|f(y)-f(b)|<\eps$. Por el límite de $g$, hay $\delta$ con $0<|x-x_0|<\delta\Rightarrow|g(x)-b|<\eta$. Encadenando, listo. $\blacksquare$
\end{demo}
\begin{trampa}{HACE FALTA CONTINUIDAD}
Sin continuidad de $f$ en $b$ falla: $\fl{\frac{\sin x}x}$ no tiene límite en $+\infty$ aunque $\frac{\sin x}x\to0$, porque $\fl\cdot$ salta en 0.
\end{trampa}

\section{Límites generalizados}
\agregado
\begin{sello}{LATERALES}
$\lim_{x\to x_0^+}f=L$: la misma definición, pero solo con $x_0<x<x_0+\delta$. (Análogo por izquierda.)\\[2pt]
\textbf{Teorema:} $\lim_{x\to x_0}f=L\iff$ existen los dos laterales y valen $L$.
\end{sello}
\begin{sello}{LÍMITES INFINITOS Y EN EL INFINITO}
$\lim_{x\to x_0}f=+\infty\iff\forall K\ \exists\delta>0:\ 0<|x-x_0|<\delta\Rightarrow f(x)>K$.\\[2pt]
$\lim_{x\to+\infty}f=L\iff\forall\eps>0\ \exists M:\ x>M\Rightarrow|f(x)-L|<\eps$.\\[2pt]
$\lim_{x\to+\infty}f=+\infty\iff\forall K\ \exists M:\ x>M\Rightarrow f(x)>K$.
\end{sello}
\begin{tip}{EL MISMO JUEGO}
Siempre es «te piden algo, vos respondés con un entorno». Cambia qué es el entorno: $(x_0-\delta,x_0+\delta)$ para un punto, $(M,+\infty)$ para el infinito; y $(L-\eps,L+\eps)$ o $(K,+\infty)$ para el valor.
\end{tip}

\section{Cálculo de límites e indeterminaciones}
\agregado
\begin{metodo}{QUÉ HACER SEGÚN LO QUE APARECE}
\begin{enumerate}
\item \textbf{Sustituir.} Si no aparece nada raro, el límite es el valor (continuidad).
\item \textbf{$\frac00$ con polinomios:} factorizar $(x-x_0)$ arriba y abajo y cancelar. Es legal: el límite solo mira $x\neq x_0$.
\item \textbf{$\frac00$ con raíces:} multiplicar por el conjugado: $(\sqrt A-\sqrt B)(\sqrt A+\sqrt B)=A-B$.
\item \textbf{Senos:} $\lim_{u\to0}\frac{\sin u}u=1$. Acomodar para que el argumento del seno aparezca abajo.
\item \textbf{Acotado por infinitésimo:} da 0.
\item \textbf{Cocientes de polinomios en $\pm\infty$:} mandan los grados (0, cociente de coeficientes principales, o $\pm\infty$).
\item \textbf{$\infty-\infty$ con raíces:} conjugado.
\item \textbf{Parte entera o valor absoluto:} laterales por separado.
\end{enumerate}
\end{metodo}
\agregado
\begin{sello}{EL LÍMITE DEL SENO}
Para $0<x<\frac\pi2$: \ $\sin x<x<\tan x$ \ (comparando áreas en la circunferencia unidad).\\[2pt]
Dividiendo por $\sin x>0$: $1<\frac x{\sin x}<\frac1{\cos x}$, o sea $\cos x<\frac{\sin x}x<1$. Como $\cos x\to1$, por sándwich: $\displaystyle\lim_{x\to0}\frac{\sin x}x=1$.
\end{sello}
\begin{center}
\begin{tikzpicture}[scale=1.8]
\fill[teal!12] (0,0)--(1,0) arc (0:40:1)--cycle;
\draw[-{Stealth}] (-0.1,0)--(1.3,0); \draw[-{Stealth}] (0,-0.1)--(0,0.95);
\draw[suave] (1,0) arc (0:70:1);
\draw[acero,thick] (0,0)--(1,0.839);
\draw[rojo,thick] ({cos(40)},0)--({cos(40)},{sin(40)});
\draw[naranja,thick] (1,0)--(1,0.839);
\node[rojo,font=\scriptsize,left] at ({cos(40)},0.3) {$\sin x$};
\node[naranja,font=\scriptsize,right] at (1,0.42) {$\tan x$};
\node[teal!70!black,font=\scriptsize] at (0.5,0.13) {$x$};
\end{tikzpicture}\\
{\ttfamily\scriptsize\color{suave} Triángulo chico $\subset$ sector $\subset$ triángulo grande: áreas $\frac{\sin x}2<\frac x2<\frac{\tan x}2$.}
\end{center}
